%% kapa: exam=1 round=26 year=2015 period=1 form=A subject=economics subject_ko=경제원론 seq=33 number=73
%% source: Q-Net 감정평가사 기출문제 (cst00309) artlSeq=5138001 file=extracted/1008/5138001/제26회_감정평가사 1차_기출/제26회_감정평가사 1차_기출/1교시_감평_A형.pdf
%% conversion: kordoc 4.24.51 -> tools/gen1.py (auto); review=proofread-sonnet; flags=none
\begin{question}{73}
단기 또는 장기 총공급곡선을 오른쪽으로 이동시키는 요인으로 옳지 않은 것은?
\begin{choices}
  \item 물적자본 증가
  \item 노동인구 증가
  \item 기술지식 진보
  \item 예상 물가수준 하락
  \item 자연실업률 상승
\end{choices}
\end{question}
\answer{5}
